\chapter{API-Grenzen, Konventionen und Fehler}
\label{api:overview}\index{API}\index{Kompatibilität}
\begingroup\raggedright
Diese Referenz beschreibt den tatsächlich vorhandenen M0--M8-Umfang und den
integrierten generischen M9-Unsicherheitsworkflow. Version und Commit stehen auf
dem Titelblatt; Paketversion, Model-Schema und Ergebnisformat sind getrennte
Verträge. Der Paketroot exportiert nur \texttt{\_\_version\_\_}; Solver werden
aus ihren Modulen importiert. Der kanonische Stabilitätsvertrag ist
\nolinkurl{Doc/29_PUBLIC_API_CONTRACT.md}, ergänzt durch die dort referenzierten
M6--M9-Verträge. Das maschinelle Inventar enthält auch Forschungsbausteine und
interne Schnittstellen. Deren Erreichbarkeit bedeutet keine SemVer-Zusage.

Die Inventarmatrix \nolinkurl{Doc/latex/api-map.json} ordnet alle Quellmodule
explizit zu: stabil, gemischt, Forschung oder intern. Reexportierte Klassen
führen auf dieselbe Definition; stabile Importpfade bleiben maßgeblich.
Incidental Imports wie \texttt{numpy}, \texttt{Path} oder \texttt{dataclass}
sind keine Toolbox-Einstiegspunkte. Unterstrich-Helfer sind Implementierungsdetails.
Öffentliche Datentypen, Konstanten, Signaturen, Felder, Konstruktoren und Methoden
werden quellenbasiert inventarisiert; direkt gefundene Raise-Stellen sind keine
vollständige transitive Exception-Liste. Hier gelten zusätzlich die manuellen
Fachverträge und die validierenden Adapter.

\section{Einheiten, Reihenfolge und Veränderlichkeit}
Direkte Python-Solver erwarten SI ohne Umrechnung. Model/TOML konvertiert nur
explizit unterstützte Einheiten. Netzkoordinaten sind Meter, Indizes nullbasiert,
externe IDs getrennte Metadaten. Aktive Knoten sind sortiert; Rohfelder verweisen
über \texttt{node\_ids} auf das Netz. Zellergebnisse folgen Eingabezellreihenfolge.
DOFs sind feldweise, innerhalb $u$ je Knoten $x,y,z$; elektrische/thermische
Felder besitzen je einen skalaren DOF. Eine \texttt{Dirichlet}-Liste muss aus
der passenden DOF-Map kommen, nicht aus ungeprüften Netzindizes.

Voigt: $xx,yy,zz,yz,xz,xy$; Dehnungen verwenden technische Schubanteile.
$\mathbf E=-\nabla\phi$, $\mathbf q=-k\nabla T$. Wärmestrom nach außen positiv.
$\mathbf D\cdot\mathbf n_{\mathrm{out}}$ ist Flächenladung; angrenzende ideale
Elektrodenladung hat umgekehrtes Vorzeichen. Harmonik verwendet komplexe
Spitzenwertphasoren und $e^{+i\omega t}$; reale/komplexe Größen nicht vermischen.

Mesh-/Material-/Ergebnisarrays werden nach den jeweiligen Klassenverträgen
kopiert und auf unveränderlichen Bytes gespeichert. Frozen-Datensätze allein
machen eingebettete externe Objekte nicht automatisch tief unveränderlich;
das Inventar nennt Dekoratoren und Typen, der Fachvertrag erläutert Besitz.
Model-Dokumente und JSON-Berichte werden als getrennte Kopien herausgegeben.
Ein Study liest Inputs bei jeder Auflösung neu. Plotobjekte enthalten mutable
Matplotlib-Objekte, die der Aufrufer schließen muss.

\section{Diagnostischer Fehlervertrag}
\texttt{FEMToolboxError} ist die fachliche Basis; \texttt{ValidationError}
erweitert \texttt{ValueError} um \texttt{field}. Shape-/Unit-/Provenance-/Mesh-
und Material-Unterklassen präzisieren Ursachen. \texttt{ModelSchemaError}
nennt \texttt{path}. Dateien können zusätzlich normale I/O-Fehler erzeugen.
\texttt{ThermalNullspaceError} und \texttt{ElectricGaugeError} unterscheiden
unverankerte Komponenten und kompatible/inkompatible Lasten. Fehlende Lagerung
wird abgelehnt; keine künstliche Stabilisierung oder automatische Gauge.
Backendfehler werden sichtbar propagiert; kein stiller Ersatz einer angeforderten
PETSc-Methode durch SciPy. Fehlendes Visualisierungs-Extra erzeugt einen
konkreten Importfehler mit Installationshinweis.

Study-Fälle können physisch fehlschlagen und persistieren diesen Status.
Schema-, Pfad-, Lock- und korruptionsbedingte Fehler sind keine erfolgreich
übersprungenen Fälle. Mengentabellen behalten fehlende Werte samt Grund;
Surrogat/UQ akzeptieren keine unmarkierte Auswahl der verbliebenen Werte.
\endgroup
